\section{The complete static base under weak decay}\label{q4:sta:section}

The causal field of Theorem~\ref{q4:adj:causal-field} satisfies the
hypotheses of Proposition~\ref{prop:static:common-base}. We verify that
application at the stated weak decay, and record the boundary coordinates
needed to recover the original hypersurface. The harmonic-coordinate
improvement and the exclusion of possible additional ends belong to the
four-dimensional classification in Section~\ref{q4:cla:section}.

Write
\[
 N=u^2-|X|_g^2,\qquad
 \mathcal P=\{N>0\}\subset\operatorname{int}\Omega,\qquad
 Z=\{N=0\}\cap\operatorname{int}\Omega.
\]
The adjoint theorem gives $N\ge0$, $N\to1$ at infinity, and a full positive
collar of $S$. Thus $Z$ is compactly contained in the original interior.
On $\mathcal P$ set
\begin{equation}\label{q4:sta:base-variables}
 \begin{aligned}
 h_b&=g+N^{-1}X^\flat\otimes X^\flat,
 &C_b&=g^{-1}-u^{-2}X\otimes X=h_b^{-1},\\
 A_b&=N^{-1}X^\flat,
 &\lambda&=\sqrt N.
 \end{aligned}
\end{equation}
Here $X^\flat$ is formed with $g$. The stationary metric is
\begin{equation}\label{q4:sta:stationary-decomposition}
 \mathbf g=-N(dt-A_b)^2+h_b.
\end{equation}

\begin{proposition}[Static base for the weak four-dimensional class]
\label{q4:sta:complete-base}
The set $\mathcal P$ is connected; denote it by $\mathcal U$. On this set
$dA_b=0$, $0<\lambda<1$, and
\begin{equation}\label{q4:sta:static-equations}
 \lambda\Ric_{h_b}=\Hess_{h_b}\lambda,
 \qquad \Delta_{h_b}\lambda=0,
 \qquad \Scal_{h_b}=0.
\end{equation}
Every approach to an interior zero of $N$ has infinite $h_b$-length.
Adjoining $S$ gives a complete manifold with smooth one-sided metric
$h_b$ and lapse $\lambda$, with
\begin{equation}\label{q4:sta:boundary-data}
 h_b|_{TS}=g|_{TS},\qquad \lambda|_S=0,
 \qquad \partial_{\nu_{h_b}}\lambda=\kappa,
 \qquad \mathrm{II}_{h_b}=0,
 \qquad \kappa=\frac1{\sqrt{2m}}.
\end{equation}
The normal $\nu_{h_b}$ points into the base exterior. The attachment is
homeomorphic to the original attachment of $S$, and it is the only
finite-distance boundary attachment. Its smooth structure is as follows:
\begin{enumerate}[label=(\roman*)]
\item If $u|_S>0$, then $N$ is an original defining function and the
base boundary structure replaces $N$ by $\lambda=\sqrt N$. The reflected
metric is smooth and even, and the signed lapse is smooth and odd.
\item If $u|_S=0$, the one-sided base metric and lapse are smooth in
the original boundary structure. Their even and odd reflections in base
normal distance are at least $C^2$.
\end{enumerate}
\end{proposition}

\begin{proof}
The original exterior is complete with $S$ included and has one
asymptotically Euclidean end with compact complement. Its boundary is
smooth, nonempty, compact and connected. Choose the exponent already
used in the adjoint construction,
\[
 1<q_0<\min\{q,2\}.
\]
The original metric has $g-\delta=O_2(r^{-q_0})$, while
Theorem~\ref{q4:adj:causal-field} proves
$(u,X)=(b^0,b)+O_2(r^{-q_0})$, with
$(b^0)^2-|b|^2=1$ and $b^0>0$. The same theorem supplies smoothness up
to $S$, $u>0$ in the interior, $u\ge|X|_g$, the first Killing initial-data
equation, and
\[
 \Ric_{\mathbf g}(\partial_t,\cdot)=0,
 \qquad \Ric_{\mathbf g}=0\quad\hbox{on }\mathcal P.
\]
Its boundary equations and lapse dichotomy give the remaining hypotheses
of Proposition~\ref{prop:static:common-base}, with the positive constant
$\kappa=1/\sqrt{2m}$.

In particular, the shared twist argument needs only the differentiated
end bounds just established. After subtracting the differential of a
function equal to $b^ix^i$ far out, its test one-form is
$\beta=A_b-df=O(r^{-q_0})$. The current
\[
 Q^{ij}=uN C_b^{ik}C_b^{j\ell}(dA_b)_{k\ell}
\]
is $O(r^{-1-q_0})$. Its cutoff term on a four-dimensional annulus is
$O(R^{2-2q_0})$, which tends to zero because $q_0>1$. The compact null-set
and horizon estimates use the smooth original fields and the boundary
collars. Thus the original decay with two metric derivatives and one
derivative of $K$ suffices for every conclusion of the shared proposition,
giving the stated static equations, connectedness, completeness and attachment.
\end{proof}

\paragraph{The attachment in original coordinates.}
The collar description fixes the original boundary structure in which
the final embedding must be smooth. If $u|_S>0$, then
$dN=2\kappa X^\flat$ on $S$. In original coordinates $(N,y^A)$, smooth
division therefore gives
\[
 X^\flat=a\,dN+N\beta_A\,dy^A,
 \qquad a(0,y)=\frac1{2\kappa}.
\]
The base coordinates are $(\lambda,y)$ with $N=\lambda^2$, and
\begin{equation}\label{q4:sta:positive-lapse-base-collar}
 \begin{split}
 (h_b)_{\lambda\lambda}&=4\lambda^2g_{NN}+4a^2,\\
 (h_b)_{\lambda A}&=2\lambda(g_{NA}+a\beta_A),\\
 (h_b)_{AB}&=g_{AB}+\lambda^2\beta_A\beta_B.
 \end{split}
\end{equation}
Every original coefficient on the right is evaluated at $(\lambda^2,y)$.
The resulting smooth reflection is exactly
$(\lambda,y)\mapsto(-\lambda,y)$: the first and third coefficients are
even and the middle coefficient is odd. At $S$ the cross terms vanish
and $(h_b)_{\lambda\lambda}=\kappa^{-2}$.

If $u|_S=0$, the original normal distance $s$ instead gives
$u=sd$, $X=s^2Y_2$, and $d|_S=\kappa$, hence
\[
 h_b=g+\frac{s^2Y_2^\flat\otimes Y_2^\flat}
                    {d^2-s^2|Y_2|_g^2},
 \qquad \lambda=s\sqrt{d^2-s^2|Y_2|_g^2}.
\]
These are smooth in the original collar. The static equation gives
$\Hess_{h_b}\lambda=0$ on $S$; in base normal coordinates the even metric
and odd lapse therefore match derivatives through order two.

The proposition retains any complete ends caused by $Z$. The next
section allows those ends throughout the conformal doubling and
zero-mass argument, and excludes them only after identifying the double
as Euclidean space.
